\documentclass[11pt,a4paper]{article}

\usepackage[utf8]{inputenc}
\usepackage[T1]{fontenc}
\usepackage[spanish]{babel}
\usepackage{amsmath,amssymb,bm}
\usepackage{physics}
\usepackage[margin=2.4cm]{geometry}
\usepackage{braket}
\usepackage{hyperref}

% \Tr, \Im, \norm ya vienen del paquete physics
\newcommand{\kk}{\bm{k}}
\newcommand{\bE}{\bm{E}}
\newcommand{\sig}{\bm{\sigma}}
\newcommand{\rhoeq}{\rho^{(0)}}

\title{QXTI: modelos y m\'etodos\\
\large Haldane, semimetal de Weyl (Orenstein), matrices densidad
perturbativas y funci\'on de Green de superficie de L\'opez-Sancho}
\author{Notas de trabajo}
\date{\today}

\begin{document}
\maketitle

\noindent
Todo en unidades at\'omicas ($\hbar=e=1$). Las energ\'ias se miden en
Hartree y los vectores de onda en $a_0^{-1}$ (Bohr$^{-1}$).

%======================================================================
\section{Modelo de Haldane}
%======================================================================
Red \emph{honeycomb} con dos subredes $A,B$, hopping a primeros vecinos
$t_1$ y a segundos vecinos con fase compleja $t_2 e^{i\phi_0}$. Es un
aislante de Chern sin campo magn\'etico neto
\cite{Haldane1988}. El Hamiltoniano $2\times2$ se escribe en la base
$\{\ket{A,\kk},\ket{B,\kk}\}$ como
\begin{equation}
  H(\kk) = B_0(\kk)\,\sigma_0 + \bm{B}(\kk)\cdot\sig,
  \qquad
  \bm{B}=(B_1,B_2,B_3),
\end{equation}
con
\begin{align}
  B_0(\kk) &= 2 t_2\cos\phi_0 \sum_{i} \cos(\kk\cdot\bm{b}_i), \\
  B_1(\kk) &= t_1 \sum_{i}\cos(\kk\cdot\bm{a}_i), \qquad
  B_2(\kk) = t_1 \sum_{i}\sin(\kk\cdot\bm{a}_i), \\
  B_3(\kk) &= M_0 - 2 t_2\sin\phi_0 \sum_{i}\sin(\kk\cdot\bm{b}_i).
\end{align}
Los $\bm{a}_i$ son los tres vectores a primeros vecinos ($A\!\to\!B$) y
los $\bm{b}_i$ los de segundos vecinos ($A\!\to\!A$). Expl\'icitamente,
con constante de red $a_0$,
\begin{equation}
  \bm{a}_i \in \Big\{ (0,a_0),\ \big(\!-\tfrac{\sqrt3}{2}a_0,-\tfrac{a_0}{2}\big),\
  \big(\tfrac{\sqrt3}{2}a_0,-\tfrac{a_0}{2}\big)\Big\},
  \quad
  \bm{b}_i \in \Big\{ (\sqrt3 a_0,0),\ \big(\!-\tfrac{\sqrt3}{2}a_0,\tfrac{3}{2}a_0\big),\
  \big(\!-\tfrac{\sqrt3}{2}a_0,-\tfrac{3}{2}a_0\big)\Big\}.
\end{equation}
El t\'ermino $B_0$ (par en $\kk$) desplaza r\'igidamente las bandas; la
masa efectiva $B_3$ contiene la masa de subred $M_0$ (rompe inversi\'on)
y el t\'ermino de Haldane $\propto t_2\sin\phi_0$ (rompe TRS). El signo
relativo de $M_0$ y $t_2\sin\phi_0$ en los dos valles $K,K'$ fija el
n\'umero de Chern $C=\pm1$ frente a $C=0$. Par\'ametros por defecto:
$t_1=2.7$~eV, $t_2=0.30$~eV, $M_0=0$, $\phi_0=\pi/2$,
$a_0=2.46$~\AA\ (grafeno realista).

%======================================================================
\section{Semimetal de Weyl de Wu/Orenstein (TaAs)}
%======================================================================
Modelo m\'inimo de cuatro nodos de Weyl que rompe inversi\'on, usado para
la respuesta SHG gigante de TaAs \cite{Wu2017}. Base $4\times4$ =
orbital $\otimes$ esp\'in, con matrices de Pauli $\sigma_i$ (orbital) y
$s_i$ (esp\'in):
\begin{equation}
  H(\kk) = t\Big[ f_x\,(\sigma_x\!\otimes\! I)
    + \sin(k_y a)\,(\sigma_y\!\otimes\! I)
    + \Delta\cos(k_y a)\,(\sigma_y\!\otimes\! s_x)
    + \sin(k_z a)\,(\sigma_z\!\otimes\! s_x)\Big],
\end{equation}
\begin{equation}
  f_x(\kk) = \cos(k_x a) + m_y\big(1-\cos(k_y a)\big) + m_z\big(1-\cos(k_z a)\big).
\end{equation}
Como $s_x$ conmuta con $H$, el Hamiltoniano es bloque-diagonal en
$s_x=\pm1$; cada bloque $2\times2$ aloja un par de nodos. En el bloque
$s_x=s$ el vector $\bm{d}$ es
\begin{equation}
  d_x = t f_x,\quad
  d_y = t\big[\sin(k_y a) + s\,\Delta\cos(k_y a)\big],\quad
  d_z = t\, s\,\sin(k_z a).
\end{equation}
\paragraph{Nodos y quiralidad.} Los ceros de $\bm{d}$ dan cuatro nodos en
$k_z=0$: $s_x=+1$ en $k_y a=-\arctan\Delta$ y $s_x=-1$ en
$k_y a=+\arctan\Delta$ (separaci\'on en $k_y$ $\Rightarrow$ ruptura de
inversi\'on $\Rightarrow$ SHG $\neq0$); $k_x$ se obtiene de $f_x=0$,
$\cos(k_x a)=-m_y(1-\cos k_y a)$. La quiralidad es el signo del
Jacobiano $\det(\partial d_i/\partial k_j)$, que en $k_z=0$ vale
\begin{equation}
  \chi = \operatorname{sgn}\!\Big[-s\,\sin(k_x a)\,
        \big(\cos(k_y a) - s\,\Delta\sin(k_y a)\big)\Big]
      = \operatorname{sgn}(k_x k_y).
\end{equation}
El factor $s$ en $d_z$ (que cambia de signo entre los dos bloques de
$s_x$) es lo que hace que la quiralidad siga el patr\'on \emph{diagonal}
$\chi=\operatorname{sgn}(k_x k_y)$: dos nodos $+1$ y dos $-1$
(Nielsen--Ninomiya). Se ha verificado num\'ericamente contra el n\'umero
de Chern por flujo de Berry (Fukui--Hatsugai--Suzuki) en una esfera
alrededor de cada nodo. Par\'ametros: $t=0.8$~eV, $m_y=1$, $m_z=5$,
$\Delta=0.5$, $a=1$.

%======================================================================
\section{Matrices densidad en el motor \emph{theory} (dominio de frecuencia)}
%======================================================================
El motor anal\'itico resuelve la ecuaci\'on de Liouville en gauge de
longitud en el dominio de frecuencia, orden a orden en el campo
\cite{Hipolito2018}. Se parte del equilibrio (matriz diagonal de
Fermi--Dirac en la base de bandas)
\begin{equation}
  \rhoeq_{mn}(\kk) = f\big(\varepsilon_n(\kk)\big)\,\delta_{mn},
  \qquad
  f(\varepsilon) = \frac{1}{e^{(\varepsilon-\mu)/T}+1},
\end{equation}
y se aplica la recursi\'on (el arm\'onico $s$ sale a la frecuencia
$s\omega$):
\begin{equation}
  \boxed{\;
  \rho^{(s)}_{mn}(s\omega) =
    \frac{\big[\bE\cdot D_{\kk}\,\rho^{(s-1)}\big]_{mn}}
         {\,s\omega - \omega_{mn} + i\gamma\,},
  \qquad \omega_{mn}=\varepsilon_m-\varepsilon_n. \;}
\end{equation}
Aqu\'i $D_{\kk}$ es el \textbf{gradiente covariante} (gauge-invariante)
\begin{equation}
  D_{\kk}\rho = \nabla_{\kk}\rho - i\,[\mathcal{A},\rho],
  \qquad
  \mathcal{A}^{\alpha}_{mn} = i\,\frac{v^{\alpha}_{mn}}{\varepsilon_m-\varepsilon_n}\ (m\neq n),\quad
  \mathcal{A}^{\alpha}_{nn}=0,
\end{equation}
con $v^{\alpha}=\partial H/\partial k_\alpha$ la velocidad y
$\mathcal{A}$ la conexi\'on de Berry inter-banda. En la pr\'actica
$\nabla_{\kk}$ se calcula por diferencias finitas alineadas por enlaces
de Wilson (transporte paralelo), lo que combina las dos piezas
($\nabla_{\kk}\rho$ y $-i[\mathcal{A},\rho]$) de forma individualmente
gauge-invariante. Expl\'icitamente, el primer orden es
\begin{equation}
  \rho^{(1)}_{mn}(\omega) =
    \frac{E_\alpha\,\mathcal{A}^{\alpha}_{mn}\,\big(f_n-f_m\big)}
         {\omega-\omega_{mn}+i\gamma}.
\end{equation}
La corriente arm\'onica y las susceptibilidades se obtienen de la traza
\begin{equation}
  J^{(s)}_{\alpha}(s\omega) = \sum_{\kk} w_{\kk}\,
      \Tr\!\big[(-v^{\alpha})\,\rho^{(s)}(\kk,s\omega)\big],
  \qquad
  \sigma^{(s)} = \frac{J^{(s)}}{\text{(campos)}},\quad
  \chi^{(s)}=\frac{\sigma^{(s)}}{(2i\,s\omega)}\ \text{(convenci\'on)} .
\end{equation}
Para orden 2 se implementa adem\'as una versi\'on \emph{grid-based}:
$\rho^{(1)}(\kk,\omega)$ se vectoriza sobre toda la malla, el gradiente
covariante se toma con \texttt{np.roll} sobre el \emph{mesh}, y se
ensambla $\rho^{(2)}(2\omega)$ simetrizado en $(j,k)$; es 50--100$\times$
m\'as r\'apido y preserva las simetr\'ias del cristal.

%======================================================================
\section{Matrices densidad perturbativas del motor CMD (dominio del tiempo)}
%======================================================================
El motor CMD resuelve la \emph{misma} f\'isica pero en el dominio del
tiempo, propagando la ecuaci\'on de Liouville--von-Neumann en gauge de
longitud orden a orden en la amplitud del campo:
\begin{equation}
  \boxed{\;
  \frac{d\rho^{(s)}_{nm}}{dt} =
    -\big(i\,\omega_{nm} + \gamma_{nm}\big)\,\rho^{(s)}_{nm}
    + \big[\bE(t)\cdot D_{\kk}\,\rho^{(s-1)}\big]_{nm}. \;}
\end{equation}
El gradiente covariante $D_{\kk}\rho=\nabla_{\kk}\rho-i[\mathcal{A},\rho]$
es el mismo que en \emph{theory} (transporte paralelo / enlaces de
Wilson); las tasas $\gamma_{nm}$ modelan la relajaci\'on ($T_1/T_2$). El
orden cero es el equilibrio diagonal de Fermi--Dirac, y para
$s=1,\dots,s_{\max}$ el t\'ermino fuente
$\mathrm{src}^{(s)}(t)=[\bE(t)\cdot D_{\kk}\rho^{(s-1)}]$ se construye del
orden anterior (parte intra-banda por $\nabla_{\kk}$ + parte inter-banda
por el conmutador de Berry).

\paragraph{Integrador.} Para cada elemento $nm$ la ODE es lineal y se
propaga con un integrador exponencial trapezoidal exacto en la parte
homog\'enea:
\begin{equation}
  \rho^{(s)}[t{+}1] = P(\Delta t)\,\rho^{(s)}[t]
    + \frac{\Delta t}{2}\Big(P(\Delta t)\,\mathrm{src}[t] + \mathrm{src}[t{+}1]\Big),
  \qquad
  P(\Delta t) = e^{-(\gamma_{nm} + i\,\omega_{nm})\,\Delta t}.
\end{equation}
La corriente total es $J_\alpha(t)=\sum_{\kk}w_{\kk}\Tr[(-v^{\alpha})\rho(t)]$
y los arm\'onicos se extraen por FFT. \emph{Relaci\'on con theory:} la
recursi\'on de frecuencia de la Sec.~3 es exactamente la transformada de
Fourier de esta recursi\'on temporal (el denominador
$s\omega-\omega_{mn}+i\gamma$ es la respuesta estacionaria de $P$ ante un
forzado $e^{-is\omega t}$). Por eso ambos motores deben coincidir, y
coinciden ($\sim1.8\%$ en $\sigma^{(1)}$).

%======================================================================
\section{LDOS de superficie: L\'opez-Sancho}
%======================================================================
El DOS de volumen se obtiene por diagonalizaci\'on directa,
\begin{equation}
  g(E) = \sum_{\kk} w_{\kk}\sum_n L_\eta\big(E-\varepsilon_n(\kk)\big)
       = -\frac{1}{\pi}\sum_{\kk} w_{\kk}\,\Im\,\Tr\big[(E+i\eta)I - H(\kk)\big]^{-1},
\end{equation}
es decir el DOS de la funci\'on de Green ensanchada. Para \emph{estados de
superficie} (p.\,ej.\ los \emph{Fermi arcs}) hace falta un sistema
semi-infinito; esto se resuelve con la \textbf{decimaci\'on iterativa de
L\'opez-Sancho} \cite{Sancho1985}.

\paragraph{Bloques de capa principal.} Se elige un eje normal a la
superficie y se hace la transformada de Fourier de $H$ a lo largo de ese
eje (aproximaci\'on de capa vecina):
\begin{equation}
  H(k_n) = H_{00} + H_{01}\,e^{i k_n c} + H_{01}^{\dagger}\,e^{-i k_n c},
  \qquad
  H_{00} = \big\langle H(k_n)\big\rangle_{k_n},\quad
  H_{01} = \big\langle H(k_n)\,e^{-i k_n c}\big\rangle_{k_n},
\end{equation}
donde $H_{00}$ es el bloque intra-capa y $H_{01}$ el acoplo entre capas
adyacentes, para cada $\kk_\parallel$ fijo.

\paragraph{Decimaci\'on.} Con $z=E+i\eta$ se inicializa
$\varepsilon_s=\varepsilon=H_{00}$, $\alpha=H_{01}$, $\beta=H_{01}^{\dagger}$
y se itera
\begin{align}
  g &= \big(zI-\varepsilon\big)^{-1}, \\
  \varepsilon_s &\leftarrow \varepsilon_s + \alpha\,g\,\beta, \\
  \varepsilon   &\leftarrow \varepsilon + \alpha\,g\,\beta + \beta\,g\,\alpha, \\
  \alpha &\leftarrow \alpha\,g\,\alpha, \qquad
  \beta \leftarrow \beta\,g\,\beta,
\end{align}
hasta que $\norm{\alpha},\norm{\beta}\to0$ (convergencia exponencial:
cada iteraci\'on duplica el n\'umero de capas efectivas). El t\'ermino
$\varepsilon_s$ acumula $\alpha g\beta$ (superficie ``inferior'') mientras
que una self-energy gemela $\varepsilon_s^{\text{top}}$ acumula
$\beta g\alpha$ (superficie ``superior''); sus estados de superficie son
complementarios (en un slab de Weyl, los \emph{Fermi arcs} conectan los
nodos de una manera en una cara y de la manera complementaria en la otra).

\paragraph{Funci\'on de Green y LDOS de superficie.}
\begin{equation}
  G_s(\kk_\parallel, z) = \big(zI - \varepsilon_s\big)^{-1},
  \qquad
  \boxed{\;
  A_{\text{surf}}(\kk_\parallel, E) =
    -\frac{1}{\pi}\,\Im\,\Tr\, G_s(\kk_\parallel, E+i\eta). \;}
\end{equation}
Barriendo $\kk_\parallel$ a energ\'ia fija $E=E_F$ se obtiene el mapa de
arcos de Fermi; barriendo $E$ a $\kk_\parallel$ fijo, la LDOS de
superficie. (En el modelo de Weyl de la Sec.~2 este mapa reproduce los
arcos que conectan nodos de quiralidad opuesta.)

%======================================================================
\begin{thebibliography}{9}
\bibitem{Haldane1988} F.\,D.\,M.\ Haldane, \emph{Model for a Quantum
  Hall Effect without Landau Levels}, Phys.\ Rev.\ Lett.\ \textbf{61},
  2015 (1988).
\bibitem{Wu2017} L.\ Wu, S.\ Patankar, T.\ Morimoto, \dots\ J.\
  Orenstein, \emph{Giant anisotropic nonlinear optical response in
  transition metal monopnictide Weyl semimetals}, Nat.\ Phys.\
  \textbf{13}, 350 (2017).
\bibitem{Hipolito2018} F.\ Hip\'olito, A.\ Taghizadeh, T.\,G.\ Pedersen,
  \emph{Nonlinear optical response of 2D materials\dots}, Phys.\ Rev.\ B
  \textbf{98}, 205420 (2018); arXiv:1802.01430.
\bibitem{Sancho1985} M.\,P.\ L\'opez Sancho, J.\,M.\ L\'opez Sancho,
  J.\ Rubio, \emph{Highly convergent schemes for the calculation of bulk
  and surface Green functions}, J.\ Phys.\ F \textbf{15}, 851 (1985).
\end{thebibliography}

\end{document}
